\chapter{คอมพิวเตอร์วิทัศน์และสวนศาสตร์เคาะเสียงจำแนกความสุกแก่ทุเรียน}
\label{chap:ch04}

\section*{แผนบริหารการสอนประจำบทที่ 4}
\addcontentsline{toc}{section}{แผนบริหารการสอนประจำบทที่ 4}

\begin{planbox}{หัวข้อเนื้อหาประจำบท}
\begin{enumerate}[topsep=2pt, itemsep=1pt, partopsep=0pt, parsep=0pt, leftmargin=1.5em]
    \item ปัญหาการตัดทุเรียนอ่อนและมาตรฐานเปอร์เซ็นต์น้ำหนักเนื้อแห้งเพื่อการส่งออก
    \item การวิเคราะห์สเปซสีและข้อจำกัดของระบบสี RGB ภายใต้สภาพแสงเงาหน้างาน
    \item การแปลงระบบสี HSV และการสกัดรอยแยกสีร่องพูของหนามทุเรียนด้วยคอมพิวเตอร์วิทัศน์
    \item ฟิสิกส์การสั่นพ้องทางสวนศาสตร์และสเปกตรัมความถี่การเคาะผลทุเรียน
    \item การแปลงฟูเรียร์อย่างเร็วและโมเดลบูรณาการพหุรูปแบบเพื่อจำแนกความสุกแก่
\end{enumerate}
\end{planbox}

\begin{planbox}{วัตถุประสงค์เชิงพฤติกรรม}
\begin{enumerate}[topsep=2pt, itemsep=1pt, partopsep=0pt, parsep=0pt, leftmargin=1.5em]
    \item ระบุเกณฑ์มาตรฐานเปอร์เซ็นต์เนื้อแห้งขั้นต่ำของทุเรียนพันธุ์หมอนทองและชะนีได้อย่างถูกต้อง
    \item อธิบายข้อแตกต่างและประโยชน์ของการแปลงสเปซสีจาก RGB ไปเป็น HSV ได้อย่างแม่นยำ
    \item เขียนโปรแกรมประมวลผลภาพเพื่อหาอัตราส่วนสีน้ำตาลอมเหลืองบริเวณร่องหนามทุเรียนได้
    \item วิเคราะห์ความถี่เรโซแนนซ์ธรรมชาติจากการเคาะผลทุเรียนด้วยอัลกอริทึม FFT ได้อย่างถูกต้อง
\end{enumerate}
\end{planbox}

\newpage

\begin{planbox}{วิธีสอนและกิจกรรมการเรียนการสอน}
\begin{enumerate}[topsep=2pt, itemsep=1pt, partopsep=0pt, parsep=0pt, leftmargin=1.5em]
    \item บรรยายหลักการเคมีกายภาพการสะสมแป้งในเนื้อทุเรียนและทฤษฎีการสั่นสะเทือนทางกล
    \item สาธิตการถ่ายภาพผลทุเรียนจริงบนสายพานลำเลียงและการตัดค่าสีผ่าน OpenCV
    \item ปฏิบัติการทดลองบันทึกเสียงเคาะทุเรียนต่างระดับความสุกด้วยไมโครโฟนความไวสูง
    \item วิเคราะห์กราฟสเปกตรัมความถี่และคำนวณเปอร์เซ็นต์เนื้อแห้งเปรียบเทียบกับการอบแห้งในเตา
\end{enumerate}
\end{planbox}

\begin{planbox}{สื่อการเรียนการสอน}
\begin{enumerate}[topsep=2pt, itemsep=1pt, partopsep=0pt, parsep=0pt, leftmargin=1.5em]
    \item เอกสารคำสอนบทที่ 4 สไลด์นำเสนอ และชุดตัวอย่างภาพถ่ายทุเรียนหลากระดับความสุก
    \item กล้องดิจิทัลอุตสาหกรรม USB และไฟส่องสว่างวงแหวน LED อุณหภูมิสีคงที่
    \item ไมโครโฟนคอนเดนเซอร์ I2S หรือ USB พร้อมไม้เคาะทุเรียนปลายยางมาตรฐาน
    \item ซอฟต์แวร์ประมวลผลสัญญาณเสียงและภาพ Python (OpenCV และ SciPy)
\end{enumerate}
\end{planbox}

\begin{planbox}{การวัดและประเมินผล}
\begin{enumerate}[topsep=2pt, itemsep=1pt, partopsep=0pt, parsep=0pt, leftmargin=1.5em]
    \item การทดสอบความรู้เรื่องการตัดแยกช่วงสี HSV และการประมวลผลความถี่เสียง
    \item การประเมินผลงานโค้ดตัดภาพร่องพูทุเรียนและการหาค่าความถี่เรโซแนนซ์
    \item การประเมินความแม่นยำในการจำแนกทุเรียนแก่และทุเรียนอ่อนจากชุดข้อมูลทดสอบ
\end{enumerate}
\end{planbox}

\clearpage

\section{ปัญหาทุเรียนอ่อนและเกณฑ์มาตรฐานน้ำหนักเนื้อแห้ง}

ปัญหา \textbf{ทุเรียนอ่อน (Premature Durian)\index{ทุเรียนอ่อน (Premature Durian)}\index{Premature Durian}} ถือเป็นวิกฤติร้ายแรงของห่วงโซ่อุปทานส่งออกผลไม้ไทย การเก็บเกี่ยวทุเรียนก่อนถึงระยะสุกแก่ทางสรีรวิทยาทำให้เนื้อทุเรียนไม่มีการเปลี่ยนรูปจากแป้งเป็นน้ำตาล เนื้อแข็งกระด้าง ไม่มีกลิ่นหอมเฉพาะตัว และไม่สามารถบ่มให้สุกได้

ตามมาตรฐานสินค้าเกษตรแห่งชาติ (มกษ. 3-2556) กำหนดเกณฑ์ \textbf{เปอร์เซ็นต์น้ำหนักเนื้อแห้งขั้นต่ำ (Dry Matter Percentage หรือ \%DM)} สำหรับทุเรียนแต่ละสายพันธุ์ดังนี้
\begin{itemize}[leftmargin=1.5em]
    \item พันธุ์หมอนทอง (*Monthong*) ต้องมีเนื้อแห้งไม่น้อยกว่าร้อยละ 32
    \item พันธุ์ชะนี (*Chanee*) ต้องมีเนื้อแห้งไม่น้อยกว่าร้อยละ 30
    \item พันธุ์กระดุมทอง (*Kradum Thong*) ต้องมีเนื้อแห้งไม่น้อยกว่าร้อยละ 27
    \item พันธุ์ก้านยาว (*Kan Yao*) ต้องมีเนื้อแห้งไม่น้อยกว่าร้อยละ 32
\end{itemize}

วิธีมาตรฐานดั้งเดิมต้องใช้การเจาะเนื้อผลไม้ไปอบแห้งในตู้อบที่อุณหภูมิ $105\,^\circ\text{C}$ เป็นเวลา $4 - 6\text{ ชั่วโมง}$ ซึ่งเป็นการทำลายผลผลิต (Destructive Testing) และไม่สามารถตรวจวัดผลผลิตได้ทุกผลบนสายพาน

\section{การวิเคราะห์ความสุกแก่ด้วยคอมพิวเตอร์วิทัศน์}

การเปลี่ยนแปลงภายนอกของผลทุเรียนเมื่อเข้าสู่ระยะสุกแก่ทางสรีรวิทยา จะเกิดการขยายตัวของพู ส่งผลให้ร่องระหว่างหนาม (Inter-thorn Grooves) แยกตัวออกจากกัน และเนื้อเยื่อบริเวณร่องพูเปลี่ยนสีจากสีเขียวเข้มไปเป็นสีน้ำตาลอมเหลืองหรือสีฟางข้าว

\subsection{ข้อจำกัดของระบบสี RGB และการแปลงสู่สเปซสี HSV}

การประมวลผลภาพบนสายพานลำเลียงจริงต้องเผชิญกับเงาตกกระทบและความผันผวนของแสง ระบบสี \textbf{RGB (Red, Green, Blue)} จะรวมระดับความสว่าง (Intensity) ไว้ในทุกช่องสัญญาณ ทำให้เมื่อผลไม้โดนเงา ค่าพิกเซลจะดรอปลงใกล้ศูนย์ คอมพิวเตอร์จึงสับสนคิดว่าเป็นจุดช้ำหรือสีดำ

การแปลงภาพไปสู่สเปซสี \textbf{HSV (Hue, Saturation, Value)\index{สเปซสีเอชเอสวี (HSV)}\index{HSV Color Space}} จะช่วยแยกข้อมูล \textbf{เฉดสีบริสุทธิ์ (Hue: H)} ออกจากระดับความเข้มของแสงสว่าง (Value: V) ทำให้ระบบตรวจจับสีสุกแก่ของร่องพูได้อย่างคงที่แม้จะมีเงาบดบัง

\begin{equation}
H = \begin{cases}
0^\circ & \text{เมื่อ } \Delta = 0 \\
60^\circ \times \left( \frac{G - B}{\Delta} \bmod 6 \right) & \text{เมื่อ } C_{\max} = R \\
60^\circ \times \left( \frac{B - R}{\Delta} + 2 \right) & \text{เมื่อ } C_{\max} = G \\
60^\circ \times \left( \frac{R - G}{\Delta} + 4 \right) & \text{เมื่อ } C_{\max} = B
\end{cases}
\label{eq:rgb_to_hsv}
\end{equation}

โดยที่ $C_{\max} = \max(R, G, B)$, $C_{\min} = \min(R, G, B)$, และ $\Delta = C_{\max} - C_{\min}$

\begin{pythonbox}[การตัดแยกสีร่องพูทุเรียนด้วย OpenCV ภาษา Python]
import cv2
import numpy as np

def analyze_durian_groove_maturity(image_path: str):
    image = cv2.imread(image_path)
    hsv = cv2.cvtColor(image, cv2.COLOR_BGR2HSV)
    
    # กำหนดช่วงสีน้ำตาลอมเหลืองของร่องพูทุเรียนแก่ในระบบ OpenCV HSV
    # H: 15-35, S: 50-220, V: 50-255
    lower_mature_groove = np.array([15, 50, 50])
    upper_mature_groove = np.array([35, 220, 255])
    
    # สร้าง Mask และกำจัดสัญญาณรบกวนด้วย Morphological Opening
    mask = cv2.inRange(hsv, lower_mature_groove, upper_mature_groove)
    kernel = cv2.getStructuringElement(cv2.MORPH_RECT, (3, 3))
    cleaned_mask = cv2.morphologyEx(mask, cv2.MORPH_OPEN, kernel)
    
    # คำนวณอัตราส่วนพื้นที่ร่องหนามแก่เทียบกับพื้นที่ผลทั้งหมด
    total_fruit_pixels = np.count_nonzero(hsv[:, :, 2] > 20)
    mature_groove_pixels = np.count_nonzero(cleaned_mask)
    groove_ratio = mature_groove_pixels / max(total_fruit_pixels, 1)
    
    return {
        "groove_ratio": groove_ratio,
        "is_visually_mature": groove_ratio > 0.18
    }
\end{pythonbox}

\section{ฟิสิกส์การสั่นพ้องทางสวนศาสตร์เคาะผลทุเรียน}

ภูมิปัญญาชาวสวนทุเรียนดั้งเดิมใช้ไม้เคาะผลทุเรียนเพื่อฟังเสียง \textbf{การสั่นพ้องทางสวนศาสตร์ (Acoustic Resonance)\index{การสั่นพ้องทางสวนศาสตร์ (Acoustic Resonance)}\index{Acoustic Resonance}} เมื่อทุเรียนสุกแก่ เนื้อภายในจะเกิดการคลายตัวและแยกตัวออกจากเปลือกด้านใน เกิดเป็นโพรงอากาศขนาดเล็ก (Air Cavity) 

ในเชิงฟิสิกส์การสั่นสะเทือน ผลทุเรียนเปรียบเสมือนระบบก้อนมวลและสปริงหลายมิติ (Multidimensional Mass-Spring System)
\begin{itemize}[leftmargin=1.5em]
    \item \textbf{ทุเรียนอ่อน:} เนื้อทุเรียนยังแน่นติดสนิทกับเปลือก ความแข็งตึง (Stiffness: $k$) มีค่าสูง มวลที่เคลื่อนที่ร่วมมีมาก ความถี่เรโซแนนซ์ธรรมชาติจะอยู่ในย่านความถี่สูง \textbf{450 ถึง 650 Hz} ทำให้ได้ยินเสียงแหลมทึบ
    \item \textbf{ทุเรียนแก่พร้อมตัด:} เนื้อทุเรียนนุ่มลงและเกิดโพรงอากาศรอบพู ความแข็งตึงสัมผัสลดลงอย่างมีนัยสำคัญ ส่งผลให้ความถี่เรโซแนนซ์ธรรมชาติลดลงมาอยู่ในย่านความถี่ต่ำ \textbf{180 ถึง 320 Hz} ทำให้ได้ยินเสียงโปร่งกลวงชัดเจน
\end{itemize}

\subsection{การวิเคราะห์สเปกตรัมความถี่ด้วยการแปลงฟูเรียร์อย่างเร็ว}

สัญญาณเสียงที่บันทึกได้จากการเคาะจะถูกนำมาแปลงจากโดเมนเวลา (Time Domain) ไปสู่โดเมนความถี่ (Frequency Domain) ด้วยอัลกอริทึม \textbf{การแปลงฟูเรียร์อย่างเร็ว (Fast Fourier Transform หรือ FFT\index{การแปลงฟูเรียร์อย่างเร็ว (FFT)}\index{Fast Fourier Transform (FFT)})}

\begin{equation}
X(k) = \sum_{n=0}^{N-1} x(n) \exp\left( -j \frac{2\pi k n}{N} \right)
\label{eq:fft}
\end{equation}

ตำแหน่งของยอดคลื่นหลัก (Dominant Peak Frequency: $f_{\text{dom}}$) สัมพันธ์เชิงเส้นกับเปอร์เซ็นต์เนื้อแห้ง (\%DM) ดังแบบจำลองการถดถอยที่ได้จากการทดลองจริงของคณะวิจัย

\begin{equation}
\% \text{DM} = 48.5 - 0.052 \times f_{\text{dom}} \quad (R^2 = 0.94)
\label{eq:acoustic_dm}
\end{equation}

\begin{pythonbox}[การวิเคราะห์ความถี่เสียงเคาะทุเรียนด้วย FFT ภาษา Python]
import numpy as np
from scipy.fft import rfft, rfftfreq

def analyze_acoustic_resonance(audio_samples: np.ndarray, sample_rate: int = 44100):
    # ปรับสัญญาณให้มีค่าเฉลี่ยเป็นศูนย์และใส่หน้าต่างแฮนนิง (Hanning Window)
    centered = audio_samples - np.mean(audio_samples)
    windowed = centered * np.hanning(len(centered))
    
    # คำนวณ FFT สเปกตรัมความถี่
    N = len(windowed)
    yf = np.abs(rfft(windowed))
    xf = rfftfreq(N, 1.0 / sample_rate)
    
    # คัดกรองเฉพาะย่านความถี่การสั่นพ้องของทุเรียน (100 - 1000 Hz)
    valid_idx = np.where((xf >= 100) & (xf <= 1000))
    dominant_freq = xf[valid_idx][np.argmax(yf[valid_idx])]
    
    # ประมาณการเปอร์เซ็นต์เนื้อแห้ง
    estimated_dm = 48.5 - (0.052 * dominant_freq)
    
    return {
        "dominant_freq_hz": dominant_freq,
        "estimated_dm_pct": estimated_dm,
        "is_mature": estimated_dm >= 32.0
    }
\end{pythonbox}

\section*{คำถามทบทวนท้ายบทที่ 4}
\addcontentsline{toc}{section}{คำถามทบทวนท้ายบทที่ 4}

\begin{enumerate}
    \item จงระบุเปอร์เซ็นต์น้ำหนักเนื้อแห้งขั้นต่ำตามมาตรฐาน มกษ. 3-2556 ของทุเรียนพันธุ์หมอนทองและชะนี พร้อมอธิบายผลเสียทางเศรษฐกิจหากเก็บเกี่ยวทุเรียนไม่ได้มาตรฐาน
    \item เหตุใดการวิเคราะห์สีร่องพูทุเรียนบนสายพานลำเลียงจึงต้องแปลงระบบสีจาก RGB เป็น HSV
    \item อธิบายหลักการทางฟิสิกส์การสั่นสะเทือนที่ทำให้ความถี่เรโซแนนซ์จากการเคาะผลทุเรียนแก่ต่ำกว่าผลทุเรียนอ่อน
    \item หากวัดความถี่เรโซแนนซ์หลักของทุเรียนหมอนทองได้ 240 Hz จงประเมินเปอร์เซ็นต์เนื้อแห้งตามแบบจำลองทางสถิติ พร้อมระบุว่าทุเรียนผลนี้ผ่านเกณฑ์ส่งออกหรือไม่
\end{enumerate}
\clearpage
